\subsection{Questão 58}

\subsubsection{Enunciado}

Em um circuito \(RLC\) série oscilante, \(R = 16{,}0\,\Omega\), 
\(C = 31{,}2\,\mu\text{F}\), \(L = 9{,}20\,\text{mH}\) e

\[
\mathcal{E} = \mathcal{E}_m \sin(\omega_d t)
\]

com \(\mathcal{E}_m = 45{,}0\,\text{V}\) e 
\(\omega_d = 3000\,\text{rad/s}\).

No instante \(t = 0{,}442\,\text{ms}\), determine:

\begin{enumerate}
    \item[(a)] a taxa \(P_G\) com a qual a energia está sendo fornecida pelo gerador;

    \item[(b)] a taxa \(P_C\) com a qual a energia do capacitor está variando;

    \item[(c)] a taxa \(P_L\) com a qual a energia do indutor está variando;

    \item[(d)] a taxa \(P_R\) com a qual a energia está sendo dissipada no resistor;

    \item[(e)] se a soma de \(P_C\), \(P_L\) e \(P_R\) é maior, menor ou igual a \(P_G\).
\end{enumerate}

\subsection*{Resposta}

Dados:

\[
R = 16,0\,\Omega,
\qquad
L = 9,20\times10^{-3}\,\text{H},
\qquad
C = 31,2\times10^{-6}\,\text{F}
\]

\[
\mathcal{E}(t)=\mathcal{E}_m\sin(\omega t),
\qquad
\mathcal{E}_m=45,0\,\text{V},
\qquad
\omega=3000\,\text{rad/s}
\]

\[
t=0,442\,\text{ms}
=0,442\times10^{-3}\,\text{s}
\]

As reatâncias são

\[
X_L=\omega L
=(3000)(9,20\times10^{-3})
=27,6\,\Omega
\]

\[
X_C=\frac{1}{\omega C}
=
\frac{1}{(3000)(31,2\times10^{-6})}
=10,68\,\Omega
\]

A impedância do circuito é

\[
Z=\sqrt{R^2+(X_L-X_C)^2}
\]

\[
Z=\sqrt{16^2+(27,6-10,68)^2}
\]

\[
Z=23,29\,\Omega
\]

O ângulo de fase vale

\[
\tan\phi=
\frac{X_L-X_C}{R}
=
\frac{16,92}{16}
\]

\[
\phi=0,813\,\text{rad}
\]

A corrente máxima é

\[
I_m=\frac{\mathcal{E}_m}{Z}
=
\frac{45}{23,29}
=
1,93\,\text{A}
\]

Calculando

\[
\omega t
=
(3000)(0,442\times10^{-3})
=
1,326\,\text{rad}
\]

a corrente instantânea é

\[
i(t)
=
I_m\sin(\omega t-\phi)
\]

\[
i
=
1,93\sin(1,326-0,813)
=
0,947\,\text{A}
\]

e a tensão do gerador no instante considerado é

\[
\mathcal{E}
=
45\sin(1,326)
=
43,67\,\text{V}
\]

\subsubsection*{(a) Potência fornecida pelo gerador}

\[
P_G
=
\mathcal{E}\,i
\]

\[
P_G
=
(43,67)(0,947)
\]

\[
\boxed{P_G=41,4\,\text{W}}
\]

\subsubsection*{(b) Taxa de variação da energia no capacitor}

A amplitude da tensão no capacitor é

\[
V_{Cm}=I_mX_C
=(1,93)(10,68)
=20,64\,\text{V}
\]

Logo,

\[
v_C
=
20,64
\sin\!\left(
1,326-0,813-\frac{\pi}{2}
\right)
=
-17,98\,\text{V}
\]

Portanto,

\[
P_C=v_C\,i
\]

\[
P_C
=
(-17,98)(0,947)
\]

\[
\boxed{P_C=-17,0\,\text{W}}
\]

\subsubsection*{(c) Taxa de variação da energia no indutor}

A amplitude da tensão no indutor é

\[
V_{Lm}
=
I_mX_L
=
(1,93)(27,6)
=
53,3\,\text{V}
\]

Assim,

\[
v_L
=
53,3
\sin\!\left(
1,326-0,813+\frac{\pi}{2}
\right)
=
46,4\,\text{V}
\]

e

\[
P_L=v_L\,i
\]

\[
P_L
=
(46,4)(0,947)
\]

\[
\boxed{P_L=43,9\,\text{W}}
\]

\subsubsection*{(d) Potência dissipada no resistor}

\[
P_R=i^2R
\]

\[
P_R=(0,947)^2(16)
\]

\[
\boxed{P_R=14,4\,\text{W}}
\]

\subsubsection*{(e) Comparação entre as potências}

\[
P_C+P_L+P_R
=
-17,0+43,9+14,4
=
41,3\,\text{W}
\]

enquanto

\[
P_G=41,4\,\text{W}
\]

A pequena diferença decorre dos arredondamentos efetuados. Portanto,

\[
\boxed{
P_G=P_C+P_L+P_R
}
\]

confirmando a conservação de energia no circuito.

\vspace{1cm}
